\documentclass[aspectratio=169,11pt]{beamer}

% ─── Clean minimal theme ────────────────────────────────────────────────────
\usetheme{default}
\usecolortheme{dove}
\usefonttheme{professionalfonts}
\setbeamertemplate{navigation symbols}{}
\setbeamertemplate{footline}[frame number]
\setbeamertemplate{frametitle}{%
  \vspace{0.4em}\insertframetitle\par
}

% Colors
\definecolor{accent}{HTML}{2171B5}
\definecolor{darktext}{HTML}{333333}
\setbeamercolor{title}{fg=accent}
\setbeamercolor{frametitle}{fg=accent}
\setbeamercolor{structure}{fg=accent}
\setbeamercolor{normal text}{fg=darktext}
\setbeamercolor{itemize item}{fg=accent}
\setbeamercolor{footline}{fg=gray}
\setbeamerfont{title}{series=\bfseries,size=\Large}
\setbeamerfont{frametitle}{series=\bfseries,size=\large}

\usepackage{graphicx}
\usepackage{booktabs}
\usepackage{amsmath}

\graphicspath{{figures/}}

\title{Diffusion Mirror Descent for MuJoCo}
\subtitle{HalfCheetah-v4 Ablation Studies}
\author{Philip Nielsen}
\date{Week of March 23rd, 2026}

\begin{document}

% ─────────────────────────────────────────────────────────────────────────────
\begin{frame}
\titlepage
\end{frame}

% ─────────────────────────────────────────────────────────────────────────────
\begin{frame}{Base Performance (10 seeds)}
\begin{center}
\includegraphics[width=0.80\linewidth,keepaspectratio]{fig1_base_curve.pdf}
\end{center}
\vspace{-0.3cm}
\begin{itemize}\setlength\itemsep{2pt}
  \item Base config: MALA-2, $\eta\!=\!16$, cosine $\beta$-schedule, guided predictor.
  \item Mean final return $\approx 13{,}730 \pm 1{,}071$ (std).
        Shaded region: 95\% CI for the mean (Student-$t$, 10 seeds).
\end{itemize}
\end{frame}

% ─────────────────────────────────────────────────────────────────────────────
\begin{frame}{Policy Learning Rate Annealing}
\begin{center}
\includegraphics[width=0.95\linewidth,keepaspectratio]{fig2_policy_lr.pdf}
\end{center}
\vspace{-0.2cm}
\begin{itemize}\setlength\itemsep{2pt}
  \item Log-linear anneal of policy LR from $3\!\times\!10^{-4}$ to target over 1M steps (5 seeds each).
  \item \textbf{LR$\,\to\,$1e-5} shows the most promise ($+750$, $p\!=\!0.058$); LR$\,\to\,$3e-5 hurts.
\end{itemize}
\end{frame}

% ─────────────────────────────────────────────────────────────────────────────
\begin{frame}{$\eta$ Annealing --- Performance \& Policy Loss}
\begin{center}
\includegraphics[width=0.95\linewidth,keepaspectratio]{fig3_tfg_annealing.pdf}
\end{center}
\vspace{-0.2cm}
\begin{itemize}\setlength\itemsep{2pt}
  \item Annealing $\eta$ from 16 causes a \textbf{sharp return drop}
        and a corresponding \textbf{policy-loss spike} (proxy for entropy).
  \item More aggressive schedules $\Rightarrow$ worse collapse
        ($\eta\!\to\!1 > \eta\!\to\!2 > \eta\!\to\!4$, all 5 seeds each).
\end{itemize}
\end{frame}

% ─────────────────────────────────────────────────────────────────────────────
\begin{frame}{Why $\eta$ Annealing Fails}
\vspace{0.3cm}
\begin{itemize}\setlength\itemsep{8pt}
  \item Ideally, $\eta \to 0$ should leave the policy roughly unchanged:
        the score network has already absorbed the $Q$-guidance signal.
  \item Instead, performance \textbf{collapses} even at moderate annealing targets.
  \item This suggests that the MALA sampler is not accurate enough---the
        learned score field does not faithfully represent the target density
        without ongoing $Q$-guidance at inference time.
  \item \textbf{Hypothesis:} more MALA corrector steps $\Rightarrow$ better sampling
        $\Rightarrow$ reduced collapse.
\end{itemize}
\end{frame}

% ─────────────────────────────────────────────────────────────────────────────
\begin{frame}{More MALA Steps Do Not Fix $\eta$ Annealing}
\begin{center}
\includegraphics[width=0.95\linewidth,keepaspectratio]{fig4_mala_tfg.pdf}
\end{center}
\vspace{-0.2cm}
\begin{itemize}\setlength\itemsep{2pt}
  \item Increasing MALA steps from 2 to 4 or 8 does \textbf{not reduce} the
        performance drop from $\eta$ annealing.
  \item The collapse pattern is nearly identical, ruling out the
        na\"ive ``insufficient corrector steps'' explanation.
\end{itemize}
\end{frame}

% ─────────────────────────────────────────────────────────────────────────────
\begin{frame}{More MALA Steps Without Annealing}
\begin{center}
\includegraphics[width=0.95\linewidth,keepaspectratio]{fig5_mala_no_anneal.pdf}
\end{center}
\vspace{-0.2cm}
\begin{itemize}\setlength\itemsep{2pt}
  \item Without $\eta$ annealing, MALA-4 and MALA-8 show a \textbf{positive trend}
        ($+630$ and $+810$ resp.)\ but neither reaches $p<0.05$ with 5 seeds.
  \item MALA-8 has wider CI due to one low-performing seed (13\,177 vs.\ group mean 14\,544).
\end{itemize}
\end{frame}

% ─────────────────────────────────────────────────────────────────────────────
\begin{frame}{Statistical Significance \& Power Analysis}
\begin{center}
\includegraphics[width=0.85\linewidth,keepaspectratio]{fig6_stats.pdf}
\end{center}
\vspace{-0.2cm}
\small
Most promising: \textbf{LR$\,\to\,$1e-5} (${\sim}$12 more seeds)
and \textbf{MALA-8} (${\sim}$20 more seeds) to reach $\alpha\!=\!0.05$ at 80\% power.
\end{frame}

% ─────────────────────────────────────────────────────────────────────────────
\begin{frame}{Self-Consistency of Mirror Descent}
\smallskip
\textbf{Axiom}\; \textcolor{accent}{(C)}\; \textbf{Compositionality:}\;
$w\!\bigl(\textstyle\prod_t p_t,\;\sum_t A_t\bigr) = \prod_t w(p_t, A_t)$

\pause
\smallskip
Set all but one factor to $(1,0)$, noting $w(1,0)=1$:
\[
  w(p,A) = w(p,0)\;\cdot\; w(1,A)
\]
\pause
Substitute back into \textcolor{accent}{(C)} and cancel $\prod p_t$:
\begin{align*}
  w(p,0)&:\quad f(p_1 p_2) = f(p_1)\,f(p_2) &&\Longrightarrow\quad f(p) = p^\alpha \\
  w(1,A)&:\quad h(A_1+A_2) = h(A_1)\,h(A_2) &&\Longrightarrow\quad h(A) = e^{\eta A}
\end{align*}
\vspace{-0.4cm}

From compositionality alone: $\;w(p,A) = p^\alpha\, e^{\eta A}$.

\pause
\medskip
\textbf{Axiom}\; \textcolor{accent}{(N)}\; \textbf{Neutrality:}\;
$w(p,\, 0) = p$\;\; (no signal $\Rightarrow$ no change)\;\; $\Longrightarrow\;\alpha = 1$.

\pause
\smallskip
\begin{center}
\fcolorbox{accent}{white}{\;\;$\displaystyle
  w(p,A) = p\,e^{\eta A}
  \qquad\Longrightarrow\qquad
  \pi'(a\!\mid\!s)\;\propto\;\pi(a\!\mid\!s)\,e^{\eta\, A(s,a)}
$\;\;}
\end{center}
\vspace{-0.15cm}
\small\textcolor{gray}{Tilting a trajectory by its total advantage = tilting each step independently. The exponential tilt is the unique update with this property.}
\end{frame}

% ─────────────────────────────────────────────────────────────────────────────
\begin{frame}{Adaptive $\eta$ via KL Budget}
\smallskip
Self-consistency $\Rightarrow$ $\eta_k$ must be a \textbf{single global scalar} per update (not state-dependent).

\pause
\medskip
\textbf{KL expansion.}\; For the tilted policy $\pi_\eta(a)\propto \pi(a)\,e^{\eta A(a)}$:
\vspace{-0.15cm}
\[
  \mathrm{KL}(\pi_\eta\,\|\,\pi)
  \;=\; \eta\,\mathbb{E}_{\pi_\eta}[A] - \log\mathbb{E}_\pi[e^{\eta A}]
  \;=\; \tfrac{\eta^2}{2}\,\mathrm{Var}_\pi(A) + O(\eta^3).
\]
\vspace{-0.3cm}

\pause
Target a fixed KL budget $\Delta$, solve $\mathrm{KL}(\pi_{\eta_k}\|\pi_k)=\Delta$:
\vspace{-0.1cm}
\[
  \boxed{\;\eta_k \;\propto\; \sqrt{\frac{2\,\Delta}{\mathrm{Var}_{\pi_k}(A_k)}}\;}
\]
\vspace{-0.3cm}

\pause
\textbf{Scale invariance.}\; $A\mapsto cA$ $\;\Rightarrow\;$ $\mathrm{Var}\mapsto c^2\mathrm{Var}$ $\;\Rightarrow\;$ $\eta\mapsto\eta/c$ $\;\Rightarrow\;$ $\eta A$ unchanged.

\pause
\medskip
\textbf{Trajectory level.}\; With $G_k(\tau)=\sum_t A_k(s_t,a_t)$, the global variance decomposes:
\vspace{-0.1cm}
\[
  \eta_k \;\propto\; \sqrt{\frac{2\,\Delta}{\mathbb{E}_{s\sim d_{\pi_k}}\!\bigl[\mathrm{Var}_{a\sim\pi_k}(A_k(s,a))\bigr]}}
\]
\vspace{-0.35cm}

\small\textcolor{gray}{Adapt to uncertainty by shrinking $\Delta$, not by making $\eta$ state-dependent.}
\end{frame}


\end{document}
